% ID: SEL-D0013-BD-E08
% ESTADO: cotejada. Segunda pasada visual completa y borrador compilado revisados.
% Fuente: Selectividad/matematicas_aplicadas/matematicas_aplicadas_P_A.pdf
% SHA-256: 0305750f6b6c3cc3edbba71a19cd7cce4276935dff5a0b7d964b02eaedeaa222
% Página: 2; bloque D; ejercicio original 8.

Se ha tomado una muestra de 16 pacientes tratados por un especialista y se
ha observado que el tiempo de espera en su consulta, en minutos, ha sido de:
\[
8\quad 9.2\quad 10\quad 8.5\quad 12\quad 9\quad 11.3\quad 7\quad
8.5\quad 8.3\quad 7.6\quad 9\quad 9.4\quad 10.5\quad 8.9\quad 6.8
\]
Supongamos que el tiempo de espera en esta consulta se distribuye según una
ley Normal de varianza 4 y media desconocida.
\begin{enumerate}[label=\alph*)]
  \item \textbf{(1.5 puntos)} Halle un intervalo de confianza al 97.5\% para
  estimar el tiempo medio de espera de los pacientes tratados por este especialista.
  \item \textbf{(1 punto)} \textquestiondown{}Cuál debería ser el tamaño mínimo
  de la muestra para asegurar, con un nivel de confianza del 90\%, que el error
  cometido sea, a lo sumo, de 0.3 minutos?
\end{enumerate}
